
Il
résulte en particulier de 1.2.3 que pour tout Module gradué libre de type fini @@M0@@, le nombre d'éléments d'une base homogène de @@M1@@ ayant un degré donné ne dépend pas de la base choisie. En particulier, les entiers @@M2@@ et @@M3@@ intervenant dans la décomposition (1.2.2), définis par les conditions @@M4@@, @@M5@@, ne dépendent pas de la décomposition.

\textbf{Proposition 1.3.} \uline{Soient} @@M6@@ \uline{un schéma quasi-compact}, @@M7@@ \uline{un} @@M8@@\uline{-Module localement libre de type fini}, @@M9@@ \uline{le fibré projectif associé}. \uline{Soit} @@M10@@ \uline{un} @@M11@@\uline{-Module de présentation finie}, \uline{plat sur} @@M12@@. \uline{Il existe un entier} @@M13@@, \uline{tel que pour} @@M14@@, \uline{on ait} :
@@M15@@

Les propriétés cherchées sont de nature locale sur @@M16@@; comme @@M17@@ est quasi-compact, on peut, en prenant le plus grand des entiers trouvés, se ramener au cas où @@M18@@ est affine.

Supposons d'abord @@M19@@ noethérien.

La propriété i) résulte alors du théorème de Serre.

La propriété ii) résulte de EGA III 7.9.10, @@M20@@ étant choisi tel que la propriété i) soit vérifiée.

D'après EGA II 3.4.5, il existe un @@M21@@-Module gradué quasi-cohérent @@M22@@, de type fini, tel que @@M23@@. Par ailleurs, @@M24@@ est TN-isomorphe à @@M25@@ d'après EGA III 2.3.1; donc il existe un entier @@M26@@ tel que @@M27@@,
